\section{Introduction}
Medical image classification has benefited substantially from deeper architectures and transfer learning \cite{litjens2017,tajbakhsh2016}, but a model suitable for a retrospective benchmark is not necessarily suitable for local inference. A clinical workstation or resource-limited site must accommodate finite memory, predictable processing time, and restrictions on transmitting sensitive images \cite{kaissis2021}. These constraints make the conventional response to an accuracy shortfall---a larger backbone or a more expensive ensemble---particularly costly. The pertinent question is how far a compact network's reduced compute remains useful as the imaging task becomes harder.

Brain MRI makes this question concrete. A four-category collection of individual images distinguishes glioma, meningioma, pituitary tumor, and no tumor \cite{nickparvar}; high performance on this collection can conceal the difficulty of grading glioma across patients and imaging sequences \cite{menze2015}. The latter task changes both the label space and the acquisition setting. Consequently, accuracy on the first benchmark must not be interpreted as an estimate of grading performance on the second. Nor does a correct slice prediction, by itself, establish a correct patient-level decision. The distinction is most consequential when a minority grade is missed despite apparently favorable aggregate accuracy.

Recent work on sub-million-parameter vision networks directs limited capacity toward diverse spatial features rather than merely shrinking established backbones \cite{kyrkou2026}. Its LS-Net architecture uses parallel atrous depthwise paths to aggregate context at different receptive fields. For MRI, this design suggests an opportunity and a potential weakness: a compact model may be fast enough for local inference while predictions for borderline slices remain sensitive to changes in apparent scale. Instead of adding parameters to every prediction, one can allocate extra inference only to those predictions for which additional views may be informative. Sensitivity to an input transform is an observed operational criterion for a gate, not proof that a particular convolution branch causes an individual error.

This work asks two connected questions. First, how far can strictly sub-1M networks approach fine-tuned, multi-million-parameter backbones on a curated brain MRI classification benchmark, and does that relationship persist on a separate glioma-grading task? Second, can architecture-informed selective inference recover some of the grading gap without expanding the learned model? We compare compact scratch-trained architectures with fine-tuned larger backbones, examine class-level errors and paired significance tests, and study patient-level as well as slice-level outcomes. We then introduce multi-scale consistency-guided test-time augmentation (MCG-TTA) for LS-Net: a confidence gate selects uncertain slices, two scale views test prediction consistency, and only disagreements receive the remaining views.

Our contributions are:
\begin{enumerate}
\item We benchmark sub-1M-parameter architectures for efficiency and cross-task generalization, spanning four-class brain tumor classification and BraTS glioma grading, and quantify where compact models lose ground to larger backbones.
\item We propose \textbf{MCG-TTA}, a parameter-free, architecture-informed test-time augmentation strategy that improves slice-level F1 by up to 1.8 percentage points at roughly 1.9$\times$ baseline latency.
\end{enumerate}
